\section{Conclusion et recommandations}\label{sec:conclusion}

\paragraph{L'approche la plus appropriée.} La forêt aléatoire, réglée par validation
croisée imbriquée groupée par article, est le meilleur modèle sur les cinq grandeurs :
$R^2$ de 0,20 (Charpy) à 0,46, RMSE de 12 à \SI{31}{\percent} sous celle de la moyenne.
Ni les modèles linéaires, avec ou sans réduction de dimension, ni l'apprentissage
semi-supervisé ne font nettement mieux (section~\ref{sec:modeles}).

\paragraph{Ce que montrent les résultats.} La difficulté principale n'est pas le choix du
modèle, mais le passage à une étude inconnue : la forêt obtient un $R^2$ de 0,70 à 0,80
quand l'article testé est déjà connu, contre 0,20 à 0,46 quand il ne l'est pas. Chaque
article a ses aciers et ses conditions d'essai. Les prédictions ramènent en outre les
soudures extrêmes vers la moyenne ; c'est pourquoi le verdict « bonne / mauvaise » tiré
des seuils de la norme ISO 2560 laisse passer la plupart des soudures mauvaises
(section~\ref{sec:classification}).

\paragraph{Recommandations, avec prudence.} Les tendances suivantes ressortent à la fois
des corrélations, des modèles et de la métallurgie :
\begin{itemize}
  \item les éléments d'alliage (Mo, Mn, V, Cr, Nb) augmentent la résistance mais
    réduisent la ductilité : une bonne soudure est un compromis, à régler selon l'usage ;
  \item un traitement thermique plus long ou plus chaud abaisse la résistance ;
  \item un fort apport de chaleur au soudage et les impuretés (S, P) vont avec une
    striction plus faible ;
  \item la ténacité se juge à la température d'usage de la pièce : une soudure
    satisfaisante à \SI{0}{\celsius} peut être fragile à \SI{-40}{\celsius}.
\end{itemize}
Ces tendances viennent de données publiées, pas d'expériences contrôlées : elles
décrivent des corrélations, pas des causes. Les aciers alliés étant presque toujours
traités, par exemple, l'effet des alliages et celui du traitement ne se séparent pas.
En l'état, le modèle peut servir à trier des recettes candidates, mais ne remplace pas
les essais.

\paragraph{Pistes.} Des données plus variées (plus d'articles, d'aciers et de procédés),
des modèles par famille d'acier, et un traitement à part des lignes Charpy fixées à
\SI{100}{\joule} et \SI{28}{\joule} sont les pistes les plus prometteuses.
